% Example M1 = 2 M2. Symbol sizes are labels, not physical body radii.
\begin{tikzpicture}[textbook diagram]
  \node[anchor=west,font=\sffamily\bfseries] at (0,4.75) {Positions and center of mass: $M_1=2M_2$};
  \coordinate (O) at (0.85,0.10);
  \coordinate (P) at (3.05,2.5);
  \coordinate (Q) at (9.35,2.5);
  \coordinate (C) at (5.15,2.5);
  \draw[vector,black!50] (O) -- +(1.05,0) node[below] {$x$};
  \draw[vector,black!50] (O) -- +(0,1.0) node[left] {$y$};
  \node[below left] at (O) {$O$};
  \draw[vector,figblue] (O) -- node[left,pos=0.68] {$\boldsymbol R_1$} (P);
  \draw[vector,figred] (O) -- node[below,pos=0.78] {$\boldsymbol R_2$} (Q);
  \draw[vector,black,dashed] (O) -- node[above,pos=0.50,sloped] {$\boldsymbol R_{\rm cm}$} (C);
  \draw[guide] (P) -- (Q);
  \fill[figblue] (P) circle (2.0pt);
  \draw[figred,line width=0.9pt,fill=white] ($(Q)+(-0.065,-0.065)$) rectangle ($(Q)+(0.065,0.065)$);
  \draw[black,line width=1pt] ($(C)+(-0.07,-0.07)$) -- ($(C)+(0.07,0.07)$);
  \draw[black,line width=1pt] ($(C)+(-0.07,0.07)$) -- ($(C)+(0.07,-0.07)$);
  \node[below right] at (P) {$M_1$};
  \node[below right] at (Q) {$M_2$};
  \node[above,yshift=1.5mm] at (C) {CM};
  \draw[vector] (3.05,3.90) -- node[above] {$\boldsymbol D=\boldsymbol R_2-\boldsymbol R_1$} (9.35,3.90);
  \draw[guide] (P) -- (3.05,4.02);
  \draw[guide] (Q) -- (9.35,4.02);
  \draw[<->] (3.05,3.12) -- node[above] {$D/3$} (5.15,3.12);
  \draw[<->] (5.15,3.12) -- node[above] {$2D/3$} (9.35,3.12);
  \node[hint,anchor=west] at (6.2,0.15) {CM is a reference point,\\not a third gravitating body.};
  \draw[black!25] (0,-0.75) -- (11.55,-0.75);
  \node[anchor=west,font=\sffamily\bfseries] at (0,-1.25) {Force and acceleration are different};
  \node[anchor=west] at (0,-2.15) {Force};
  \draw[vector,figblue] (3.1,-2.15) -- (4.65,-2.15) node[above,pos=0.48] {$\boldsymbol F_1$};
  \draw[vector,figred,dashed] (9.4,-2.15) -- (7.85,-2.15) node[above,pos=0.5] {$\boldsymbol F_2$};
  \node at (6.2,-2.15) {$|F_1|=|F_2|$};
  \node[anchor=west] at (0,-3.20) {Acceleration};
  \draw[vector,figblue] (3.1,-3.20) -- (4.15,-3.20) node[above,pos=0.5] {$\boldsymbol A_1$};
  \draw[vector,figred,dashed] (9.4,-3.20) -- (7.30,-3.20) node[above,pos=0.5] {$\boldsymbol A_2$};
  \node at (5.85,-3.20) {$|A_2|=2|A_1|$};
\end{tikzpicture}
